% Generated draft: requires mathematical and source review.
\documentclass[11pt]{article}
\usepackage[margin=25mm]{geometry}
\usepackage{fontspec}
\setmainfont{Latin Modern Roman}
\usepackage{amsmath,amssymb,mathtools,unicode-math}
\setmathfont{Latin Modern Math}
\usepackage{enumitem,needspace,xurl,hyperref}
\setlength{\parindent}{0pt}
\setlength{\parskip}{5pt}
\setlength{\emergencystretch}{3em}
\newcommand{\N}{\mathbb{N}}
\newcommand{\Z}{\mathbb{Z}}
\newcommand{\Q}{\mathbb{Q}}
\newcommand{\R}{\mathbb{R}}
\newcommand{\C}{\mathbb{C}}
\newcommand{\F}{\mathbb{F}}
\newcommand{\A}{\mathbb{A}}
\newcommand{\PP}{\mathbb{P}}
\newcommand{\E}{\mathbb{E}}
\newcommand{\eps}{\varepsilon}
\newcommand{\dd}{\,\mathrm{d}}
\begin{document}
\section*{UnsolvedMath ID: 11000307}
\textbf{Source problem code:} catalogue entry
\subsection{Definitions and mathematical statement}
Problem 3.4 asks to prove that every secondary class
\[
 d_i\quad(i\ge2)
\]
is nonzero in its source-defined target. The construction, degrees, coefficients, groups/spaces and genus range of these classes are absent from the extract; nonvanishing cannot be assigned to an invented cohomology group or stable range.

The context then asks about the first class \(d_1\), viewed there as a numerical invariant on \(K_g\). Conventionally \(K_g\) is the Johnson kernel, but its precise definition and \(d_1\) normalization need source confirmation. If a separating simple closed curve \(C\subset\Sigma_g\) cuts the closed surface into genera \(h\) and \(g-h\), the source reports, for its Dehn twist \(\tau_C\in K_g\),
\[
 d_1(\tau_C)=c_g\,h(g-h)
\]
under the usual multiplicative reading of ``up to a constant depending on \(g\).'' The value of \(c_g\), twist direction and coefficient convention are not supplied and are not set to one by default.

The further task is to give an algorithm evaluating \(d_1(\phi)\) for a specified \(\phi\in K_g\), for example by analyzing its action on \(\pi_1(\Sigma_g)\). The representation of \(\phi\), based versus outer action, and output normalization must be fixed. The source's statement that no such algorithm is currently known is its historical assessment, not a present-day verification. The special separating-twist formula is not asserted to supply a general algorithm or to prove nontriviality of the higher secondary classes.
\subsection{Short English statement}
Prove nonvanishing of all the specified secondary classes beyond \(d_1\), and separately compute \(d_1\) effectively on arbitrary elements of the source's Johnson kernel. The known separating-twist formula retains its unspecified genus-dependent normalization.
\subsection{Sources}
\begin{itemize}
\item Local UnsolvedMath registry, permanent ID \(11000307\), exact source code \texttt{catalogue entry}, Problem 3.4, PDF page 363. No individual proposer or proposed year is supplied.
\item Farb, edited volume \emph{Problems on Mapping Class Groups and Related Topics}: \url{https://www.math.uchicago.edu/~farb/papers/mcgbook.pdf}. The volume has 329 explicitly typeset numbered items across chapters, repeated local numbering distinguished by PDF page.
\item Extraction pdf-font-structure, accessed 2026-07-29. Source evidence NEEDS\_REVIEW presents an open problem and requires a current primary-literature audit. Supplied OPEN-TRIAGE summary mentions a Torelli/Morita--Hain secondary-class framework and partial nontriviality, but gives no precise cases or complete references. No separate defining citation for \(d_i\) is supplied.
\item Dataset attribution: Ulam AI / UnsolvedMath Contributors, CC BY 4.0. Source rights NEEDS\_REVIEW: publicly accessible author/publisher copy, redistribution terms requiring review.
\end{itemize}
\end{document}
